\chapter*{Introduction}\label{ch:intro}
\addcontentsline{toc}{chapter}{Introduction}
\markboth{Introduction}{Introduction}

Ask double machine learning how much eligibility for a 401(k) plan raised the savings of American households in 1991 and it returns one number: about nine thousand dollars of net financial assets, with a 95 percent interval about five thousand dollars wide. Ask what the effect was at the median of the wealth distribution and the machine must be rebuilt, with a new score, a new derivation and a new run. Ask about the poorest fifth and the fitted nuisances can be reused, but the group gets its own average and its own interval. Ask how likely it is that the effect exceeds five thousand dollars and a confidence interval has no answer to give. Each answer DML does give is valid under its assumptions, and each is built separately.

This book argues for a different organization of the work. Write one probability model for outcomes given treatment and covariates, compute its posterior, and read every causal quantity off the same draws. The average treatment effect then needs no estimator of its own: it is an accounting identity, an average of fitted contrasts over a population the analyst chooses. Figure~\ref{fig:intro} shows what this buys on the 401(k) data. One quantile network is refitted 200 times under random weights (the weighted Bayesian bootstrap of Section~\ref{sec:deephd}), and each fit predicts every household's wealth with and without eligibility. Averaging the predicted contrasts over all households gives the average effect; comparing quantiles of the two predicted wealth distributions gives the effect at each quantile; averaging the contrasts within income groups gives effects from about three thousand dollars for the poorest fifth of households to eighteen thousand for the richest. All three panels come from the same 200 fits, and their spread across the fits is an approximate posterior.

\begin{figure}[ht]
\centering
\includegraphics[width=\linewidth]{intro_onefit.pdf}
\caption{One model, many answers: 9{,}915 households in the 1991 SIPP, with the nine covariates of Section~\ref{sec:deephd}. Left: the two partially linear DML estimates of the effect of 401(k) eligibility on net financial assets (random forest and lasso nuisances) and their intervals (Table~\ref{tab:k401}), beside the approximate posterior of the average effect from 200 weighted Bayesian bootstrap refits of one quantile network. The partially linear estimates weight household effects by the conditional variance of eligibility given the covariates, and the network weights them equally; Section~\ref{sec:deephd} gives interactive DML estimates of the equally weighted average. Centre: the effect at every quantile of net financial assets. Right: the mean effect by income quintile. Centre and right come from the same draws as the posterior on the left, with bands of the mean plus or minus 1.96 standard deviations; in a simulation built to resemble these data, such a band covered the true average effect in 92 percent of 1,000 data sets drawn from a network of the same kind, and in 15 percent of 200 drawn from a truth outside that family (Section~\ref{sec:deephd}).}
\label{fig:intro}
\end{figure}

The causal reading of Figure~\ref{fig:intro} assumes that, given the nine recorded covariates, eligibility is as good as randomly assigned, with overlap, consistency and no interference between households (Section~\ref{sec:ignorability}). The centre panel then compares quantiles of the two potential asset distributions over the same households. It does not give the gain of the household at a particular wealth rank, which needs the further assumptions of Section~\ref{sec:ot}.

Neyman's potential outcomes and Rubin's Bayesian imputation already contain this idea. Three obstacles kept it from practice.

\paragraph{High dimensions.} With many covariates a prior must shrink, and a prior that shrinks the wrong coefficients shrinks away the confounding. The Robins and Ritov counterexample, long read as a proof that Bayesian causal inference fails in high dimensions, shows the same leak in another setting: independent priors on many sparsely observed cells that ignore how units were selected. We show that the failure belongs to the prior, not to Bayes. The horseshoe lets large effects through and shrinks noise hard. In our designs with five strong confounders, it leaves much less regularization induced confounding than the lasso. Coupling the outcome prior to a treatment regression fitted on the whole sample reduces it further with 50 and 150 covariates, though not with 300 (Chapters~\ref{ch:ipw} and~\ref{ch:horseshoe}). Shared scales also protect instruments, covariates that affect the outcome only through treatment, and leaving them unshrunk in the outcome regression inflates variance and, with residual confounding, amplifies bias (Section~\ref{sec:sharedsim}).

In our simulation of Section~\ref{sec:ricsim}, with 200 observations and 50 covariates the coupled posterior matches least squares, which needs no regularization there, and trails DML slightly on the average effect (root mean squared error 0.084 against DML's 0.079); the 95 percent intervals of DML and of the posterior both cover the truth in about 92 percent of 500 data sets, and the posterior also gives every other summary. With 150 and with 300 covariates among the same 200 observations and the same five confounders DML is the more accurate (root mean squared error 0.096 against DML's 0.080, and 0.124 against 0.084), and the coupling gains nothing in accuracy over the independent horseshoe that the data can resolve. With thirty weak confounders among 150 covariates the order reverses: the coupled posterior is the more accurate (root mean squared error 0.113 against DML's 0.126), and its intervals cover the truth in 94 percent of the data sets against DML's 76 percent. With 300 observations and 400 covariates the comparison with a reduced form posterior (with horseshoe priors on the coefficients of the treatment and marginal outcome regressions) is mixed: DML with a cross validated lasso has the smaller bias when confounding is sparse and overlap is good, and when it is dense no procedure covers (Chapter~\ref{ch:critique}).

\paragraph{Computation.} A posterior takes thousands of Markov chain iterations where DML takes a handful of regressions, and for the horseshoe models of this book the difference is seconds: the Gibbs sampler fits the 401(k) data, 9,915 households and 56 columns, in under a second on one core (Chapter~\ref{ch:computation}).

Where the likelihood cannot be evaluated cheaply, generative Bayesian computation turns posterior computation into supervised learning (Polson and Sokolov 2023, 2026a; Nareklishvili, Polson and Sokolov 2023a). A quantile neural network trained on simulations from the model maps uniform noise and a summary of the data to draws from the posterior given that summary, and each new data set from the same design costs a forward pass per draw. The training is the expense, more than an hour per network in our horseshoe experiment, so it pays only over many data sets from one design, and for the horseshoe the sampler won. 

Chapter~\ref{ch:gente} fits the network to the observed outcomes for distributional effects. Deep partial least squares projects many covariates onto a few directions first; with 200 irrelevant covariates added to the 401(k) data it helped, but it did not recover the nine covariate answer (Section~\ref{sec:deephd}).

\paragraph{Identification.} The probability that treatment caused a particular outcome depends on how a unit's two potential outcomes are coupled and cannot be identified from their margins, which is all trials and observational studies see. Report a posterior inside the Tian and Pearl bounds, separating what more data could narrow from what only the prior fixes (Chapter~\ref{ch:tianpearl}).

The payoff of working with one posterior is decisions. A posterior gives the probability that a programme helps, the action that maximizes expected utility, and risk sensitive alternatives for a planner who cares about the worst off, all from the same draws (Chapter~\ref{ch:cdt}). We work one decision in full: whether a merger raises prices, estimated with a Bayesian version of the BLP demand model in which the conduct of firms is itself uncertain (Chapter~\ref{ch:blp}).


